\tituloI{METODOLOGÍA}

\tituloII{Método, tipo y alcance de la investigación}

\tituloIII{Método de investigación}
\parrafoIII{
    Se empleará el \textbf{método experimental y el hipotético-deductivo}. El método experimental permitirá manipular las entradas (prompts y algoritmos) para observar y medir las variaciones en las métricas de rendimiento computacional. Complementariamente, el método deductivo se aplicará para contrastar los resultados obtenidos frente a las leyes de la complejidad asintótica ($O$) y las teorías establecidas en las ciencias de la computación \cite{hernandez2014}.
}

\tituloIII{Tipo de investigación}
\parrafoIII{
    La investigación es de \textbf{tipo aplicada y tecnológica}. Es aplicada debido a que busca utilizar los conocimientos de ingeniería para evaluar la viabilidad técnica del uso de código generado por IA en entornos reales de desarrollo. Es tecnológica pues propone un marco de evaluación basado en herramientas de perfilado (profiling) para mejorar la calidad del software producido en la disciplina \cite{arias2012}.
}

\tituloIII{Alcance de la investigación}
\parrafoIII{
    El estudio posee un \textbf{alcance descriptivo-comparativo}. Su propósito es describir con precisión las métricas de tiempo y memoria de los algoritmos generados y, simultáneamente, comparar el desempeño entre distintos modelos de inteligencia artificial y las implementaciones óptimas de referencia \cite{hernandez2014}.
}

\tituloII{Materiales y métodos}

\tituloIII{Materiales}
\parrafoIII{
    Para la ejecución de las pruebas de benchmarking y el procesamiento de métricas, se utilizarán los siguientes recursos:
}
\begin{listaIII}
    \item \textbf{Hardware:} Computadora personal con arquitectura x86\_64 (procesador multinúcleo), utilizada como entorno de ejecución "Bare Metal" para evitar la latencia de virtualización.
    \item \textbf{Software de Base:} Sistema operativo Arch Linux, seleccionado por su mínima sobrecarga de procesos en segundo plano, permitiendo una medición más exacta de los recursos de CPU y RAM.
    \item \textbf{Entorno de Programación:} Intérprete de Python 3.x y compilador GCC, según el lenguaje en que se generen los algoritmos de prueba.
    \item \textbf{Instrumentos de Medición:} Herramienta \textit{Hyperfine} para la medición estadística de tiempos de ejecución y la utilidad \textit{GNU time} (o Valgrind) para el seguimiento del uso máximo de memoria residente (Peak RSS).
\end{listaIII}

\tituloIII{Métodos}
\parrafoIII{
    El procedimiento experimental se estructurará en las siguientes fases:
}
\begin{listaIII}
    \item \textbf{1) Fase de Diseño Experimental:} Selección de algoritmos de ordenación (ej. QuickSort) y grafos (ej. Dijkstra) y redacción de los prompts de prueba.
    \item \textbf{2) Fase de Generación y Limpieza:} Obtención de los scripts de código desde los modelos de IA y revisión manual mínima para asegurar la compilación/ejecución sin alterar la lógica algorítmica.
    \item \textbf{3) Fase de Benchmarking (Ejecución):} Ejecución automatizada de los algoritmos bajo diferentes cargas de datos ($n=10^2, 10^4, 10^6$) utilizando \textit{Hyperfine} para obtener medias y desviaciones estándar.
    \item \textbf{4) Fase de Análisis:} Tabulación de resultados y generación de gráficas de complejidad temporal y espacial para su interpretación técnica.
\end{listaIII}